\documentclass[final]{beamer}

% ====================
% Packages
% ====================

\usepackage{fontspec}
\usepackage[T1]{fontenc}
\usepackage{lmodern}
\usepackage[size=a1]{beamerposter}
\usepackage{./University_of_Michigan__UMich__poster_template__1_/beamerthemegemini}
\usepackage{./University_of_Michigan__UMich__poster_template__1_/beamercolorthemegemini}
\usepackage{./University_of_Michigan__UMich__poster_template__1_/beamercolorthememsu}
\usepackage{graphicx}
\usepackage{booktabs}
\usepackage{tikz}
\usepackage{pgfplots}
\pgfplotsset{compat=1.14}
\usepackage{anyfontsize}
\usepackage{hyperref}

% ====================
% Lengths
% ====================

% If you have N columns, choose \sepwidth and \colwidth such that
% (N+1)*\sepwidth + N*\colwidth = \paperwidth
\newlength{\sepwidth}
\newlength{\colwidth}
\setlength{\sepwidth}{0.025\paperwidth}
\setlength{\colwidth}{0.3\paperwidth}

\newcommand{\separatorcolumn}{\begin{column}{\sepwidth}\end{column}}

% ====================
% Title
% ====================

\title{Fused3S: GPU Kernel Fusion for Sparse Graph Transformer Attention}

\author{Project Supervisor - Prof. Sameer Kulkarni}

\institute[shortinst]{Fused Sparse-Dense Graph Attention on Tensor Cores}

% ====================
% Footer (optional)
% ====================

\footercontent{
  \href{https://github.com/fatalhawk/fused3S}{Mayank Choudhary 24110199} \hfill
  CS299 \hfill
  IIT Gandhinagar}

% ====================
% Logo (optional)
% ====================

\logoright{\includegraphics[height=5cm]{University_of_Michigan__UMich__poster_template__1_/logos/Final_IITGN-Logo-symmetric-White.png}}

% ====================
% Body
% ====================

\begin{document}

\begin{frame}[t]
\begin{columns}[t]
\separatorcolumn

% ============ FIRST COLUMN ============

\begin{column}{\colwidth}

  % -------- Overview Block --------
  \begin{block}{Overview \& Motivation}

    Graph transformers rely on sparse attention mechanisms to capture structural relationships in graph-structured data. However, the standard three-kernel pipeline incurs significant memory traffic:
    
    \begin{itemize}
      \item \textbf{SDDMM} (Sampled Dense-Dense Matrix Multiply)
      \item \textbf{Softmax} normalization
      \item \textbf{SpMM} (Sparse Matrix-Matrix Multiply)
    \end{itemize}

    \vspace{0.5em}
    
    \textbf{Fused3S} eliminates intermediate materializations by fusing all three stages into a single GPU kernel using NVIDIA Tensor Cores (WMMA), operating on block-sparse graph adjacency representations.

  \end{block}

  % -------- Key Concepts Block --------
  \begin{block}{Key Concepts}

    \begin{tabular}{ll}
      \toprule
      \textbf{Concept} & \textbf{Description} \\
      \midrule
      Block-Sparse Format & CSR $\rightarrow$ 16×16 block tiles \\
      Fused Pipeline & SDDMM $\rightarrow$ Softmax $\rightarrow$ SpMM \\
      Tensor Cores & WMMA on 16×16×16 FP16 tiles \\
      Online Softmax & Streaming algorithm for normalization \\
      Query Caching & Reused across column blocks \\
      Load Balancing & Sorted window distribution \\
      \bottomrule
    \end{tabular}

  \end{block}

  % -------- Architecture Block --------
  \begin{block}{Kernel Architecture}

    \begin{description}
      \item[\textbf{v0}:] Per-Block Softmax + 1D Layout
      
      Baseline fused kernel with per-warp softmax. Simple and functional but softmax correctness limited to individual blocks.

      \item[\textbf{v1}:] Online Softmax + 2D Layout
      
      Introduces streaming softmax maintaining global max/sum state across column blocks. Uses 2D thread layout for better organization.

      \item[\textbf{v2}:] Register-Resident Output
      
      Moves output accumulator into WMMA register fragments, eliminates per-iteration shared memory overhead, uses fragment rescaling.
    \end{description}

    \begin{figure}
      \centering
      \includegraphics[width=0.95\textwidth]{fused3s/sc_csvs/gpu-arch.png}
      \vspace{-0.5em}
      \caption*{GPU architecture for the fused kernel design}
    \end{figure}

  \end{block}

\end{column}

\separatorcolumn

% ============ SECOND COLUMN ============

\begin{column}{\colwidth}

  % -------- Roofline Analysis --------
  \begin{block}{Performance Characterization: Roofline Model}

    \begin{figure}
      \centering
      \includegraphics[width=0.95\textwidth]{fused3s/sc_csvs/roofline.png}
      \caption{Floating Point Roofline with the following points: blue = v2, orange = v1, purple = v0, green = baseline.}
    \end{figure}

    The roofline plot reveals:
    \begin{itemize}
      \item \textbf{Blue (v2):} highest arithmetic intensity and best achieved throughput among the fused kernels.
      \item \textbf{Orange (v1):} mid-tier improvement from online softmax and 2D layout.
      \item \textbf{Purple (v0):} early fused approach with correct block-level performance but lower efficiency.
      \item \textbf{Green (baseline):} naive GPU baseline with low utilization and memory-bound performance.
      \item \textbf{Trend:} v0 $\rightarrow$ v1 $\rightarrow$ v2 moves upward and rightward toward the compute roof, showing increasing optimization.
    \end{itemize}

  \end{block}

  % -------- Performance Metrics Block --------
  \begin{block}{Performance Metrics (Cora Dataset)}

    \begin{figure}
      \centering
      \includegraphics[width=0.95\textwidth]{fused3s/sc_csvs/gpu_throughput.png}
      \caption{GPU throughput comparison across baseline, v0, v1, and v2 kernels.}
    \end{figure}

    Key observations from profiling:
    
    \begin{itemize}
      \item \textbf{GPU Throughput:} v2 achieves the highest sustained throughput, clearly outperforming baseline, v0, and v1.
      \item \textbf{Baseline vs. v0:} Baseline remains strongly memory-bound, while v0 improves compute efficiency by fusing SDDMM, softmax, and SpMM.
      \item \textbf{v1 improvement:} The addition of online softmax and a 2D thread layout increases achieved throughput and raises the kernel's position on the roofline.
      \item \textbf{v2 advantage:} Register-resident output and fragment rescaling deliver the largest gain, with v2 showing the best balance of bandwidth and compute utilization.
      \item \textbf{Memory Bandwidth:} 
      \begin{itemize}
        \item Baseline: $\sim 9.96\%$ of theoretical peak
        \item v2: $\sim 37.19\%$ of theoretical peak (3.7× gain)
      \end{itemize}
      \item \textbf{Compute Throughput:}
      \begin{itemize}
        \item Baseline: $\sim 4.56\%$ of theoretical peak
        \item v2: $\sim 21.69\%$ of theoretical peak (4.8× gain)
      \end{itemize}
      \item \textbf{Relative progression:} 
      \begin{itemize}
        \item Baseline is the lowest-performing reference point.
        \item v0 gives the first fused speedup.
        \item v1 brings higher effective throughput.
        \item v2 delivers the strongest final performance and is the best candidate for production use.
      \end{itemize}
    \end{itemize}

  \end{block}

\end{column}

\separatorcolumn

% ============ THIRD COLUMN ============

\begin{column}{\colwidth}

  % -------- Implementation Details --------
  \begin{block}{Implementation Details}

    \textbf{Technology Stack:}
    \begin{itemize}
      \item CUDA Toolkit $\geq\!11.0$
      \item GPU: Compute Capability $\geq\!$sm\_80 (RTX 4050)
      \item C++17 with CMake 3.24+
    \end{itemize}

    \vspace{0.5em}
    
    \textbf{Data Processing:}
    \begin{itemize}
      \item Datasets: Cora, Citeseer
      \item Format: CSR (Compressed Sparse Row)
      \item Conversion Tools: PyTorch Geometric
    \end{itemize}

    \vspace{0.5em}
    
    \textbf{Profiling:}
    \begin{itemize}
      \item NVIDIA Nsight Compute (NCU)
      \item Baseline CPU validation
      \item Kernel correctness tests
    \end{itemize}

  \end{block}

  % -------- Key Optimizations --------
  \begin{block}{Key Optimizations}

    \begin{enumerate}
      \item \textbf{Kernel Fusion:} Eliminate intermediate matrix materializations in global memory
      
      \item \textbf{Shared Memory Pipeline:} Data flows through shared memory / registers instead of DRAM
      
      \item \textbf{Tensor Core Utilization:} WMMA instructions for dense matrix blocks within sparse structure
      
      \item \textbf{Register Blocking:} v2 keeps output accumulator in register fragments across iterations
      
      \item \textbf{Load Balancing:} Sort row windows by weight to minimize SM stalls
      
      \item \textbf{Online Softmax:} Numerically correct attention weights without separate normalization pass
    \end{enumerate}

  \end{block}

  % -------- Results & Contributions --------
  \begin{block}{Results \& Contributions}

    \begin{itemize}
      \item[$\checkmark$] Successfully implemented progressive kernel optimizations (v0, v1, v2)
      
      \item[$\checkmark$] Achieved \textbf{3-5× memory throughput improvement} over baseline
      
      \item[$\checkmark$] Demonstrated correctness via CPU baseline validation
      
      \item[$\checkmark$] Comprehensive profiling analysis with Roofline model
      
      \item[$\checkmark$] Modular design supporting multiple datasets and configurations
      
      \item[$\checkmark$] Open-source implementation with full documentation
    \end{itemize}

  \end{block}

  \begin{block}{References and Future Work}

    \textbf{Reference:}
    \begin{itemize}
      \item Z. Li and A. Chandramowlishwaran, "Fused3S: Fast Sparse Attention on Tensor Cores," May 2025, arXiv:2505.08098. [Online]. Available: https://arxiv.org/abs/2505.08098
    \end{itemize}

    \textbf{Future Work:}
    \begin{itemize}
      \item Target \textbf{> 60\% tensor core utilization} with advanced scheduling, data packing, and kernel tiling strategies.
      \item Explore finer-grained SM workload balancing, latency hiding, and register pressure reduction for further GPU performance study.
      \item Extend the fused sparse attention pipeline to larger, real-world graph workloads and mixed-precision execution.
    \end{itemize}

  \end{block}

\end{column}

\separatorcolumn

\end{columns}

\end{frame}

\end{document}
